\section{Hardness Reductions, Approximation Barrier, and Boundary Cases}
\label{app:dict-hardness}

This appendix expands Section~\ref{sec:vac-minimization}. The main text isolates a deliberately narrow algorithmic question: once a finite pool of verifier-visible candidate cells has \emph{already} been exposed, how hard is it to choose a smallest accepted subcover? The answer is that unrestricted explicit-dictionary selection already contains the full combinatorial hardness of \textsc{Set Cover}. Here we give the full reductions, record optimum preservation on both the exact and threshold sides, transfer the classical logarithmic approximation barrier, and make the boundary cases explicit.

The scope is intentionally precise. We do \emph{not} claim here that unrestricted semantic VAC over an arbitrary continuous verifier language is itself presented as an NP-complete decision problem. Instead, we isolate the more robust statement that the finite-dictionary decision problems are already NP-complete, and that the corresponding optimization problems are already hard, even when every local validity condition can be re-checked by direct pointwise inspection.

Throughout, all numerical inputs are rational and encoded in binary. In particular, the point domain
\[
U=\{u_1,\dots,u_N\}\subseteq \mathbb Q^n
\]
and the value table
\[
\{g(u_i)\}_{i=1}^N\subseteq \mathbb Q
\]
are explicit parts of the input. We write
\[
\operatorname{OPT}_{\mathrm{SC}}(\mathcal U,\mathcal S)
\]
for the optimum of a set-cover instance with universe $\mathcal U$ and family $\mathcal S$.

\subsection{Finite-Dictionary Problems and Explicit Verification}

We first restate the two dictionary-restricted decision problems together with their associated optimization problems in fully explicit form.

\paragraph{Exact-side dictionary problem.}
An instance of \textsc{Exact-VAC-Dict} consists of:
\begin{enumerate}
    \item a finite rational point set
    \[
    U=\{u_1,\dots,u_N\}\subseteq \mathbb Q^n;
    \]
    \item an explicit value table $\{g(u_i)\}_{i=1}^N$;
    \item a finite dictionary
    \[
    \mathcal D_{=} = \{\eta_1,\dots,\eta_m\},
    \qquad
    \eta_j=(A_j,a_j),
    \]
    where $A_j\subseteq U$ is explicitly listed and $a_j$ is an affine function with rational coefficients;
    \item a budget $k\in\mathbb N$.
\end{enumerate}
A dictionary entry $\eta_j=(A_j,a_j)$ is \emph{valid} if
\[
g(u)=a_j(u)
\qquad \forall u\in A_j.
\]
The decision problem asks whether there exists an index set
\[
J\subseteq \{1,\dots,m\},
\qquad
|J|\le k,
\]
such that every selected entry is valid and
\[
U\subseteq \bigcup_{j\in J} A_j.
\]
The associated optimization problem minimizes the number of selected valid entries and is denoted
\[
\operatorname{OPT}_{=}(U,\mathcal D_{=}).
\]

\paragraph{Threshold-side dictionary problem.}
An instance of \textsc{Thresh-VAC-Dict} consists of the same finite rational point set $U$, the same explicit value table, a rational threshold $t$, a finite dictionary
\[
\mathcal D_{\ge}=\{A_1,\dots,A_m\},
\qquad
A_j\subseteq U,
\]
and a budget $k\in\mathbb N$. A dictionary entry $A_j$ is \emph{valid} if
\[
g(u)\ge t
\qquad \forall u\in A_j.
\]
The decision problem asks whether there exists
\[
J\subseteq \{1,\dots,m\},
\qquad
|J|\le k,
\]
such that every selected entry is valid and
\[
U\subseteq \bigcup_{j\in J} A_j.
\]
Its optimization value is written
\[
\operatorname{OPT}_{\ge}(U,\mathcal D_{\ge},t).
\]

These two problems are the cleanest dictionary-restricted versions of verifier-aligned cover selection. All candidate local proof units are already explicit. The remaining tasks are the decision question of whether a budget-$k$ accepted subcover exists and the associated optimization question of choosing a smallest accepted subcover.

\paragraph{Proposition G.1 (Both dictionary problems belong to $\mathrm{NP}$).}
\textsc{Exact-VAC-Dict} and \textsc{Thresh-VAC-Dict} are both in $\mathrm{NP}$.

\paragraph{Proof.}
For \textsc{Exact-VAC-Dict}, a witness is the chosen index set
\[
J\subseteq \{1,\dots,m\}.
\]
The verifier checks in polynomial time that $|J|\le k$, that for every selected entry $\eta_j=(A_j,a_j)$ and every $u\in A_j$ one has
\[
g(u)=a_j(u),
\]
and that the selected subsets cover all of $U$. Since the point set, the value table, and the affine coefficients are explicit rational data, every arithmetic comparison is polynomial-time in the total input length.

For \textsc{Thresh-VAC-Dict}, the same witness format works. The verifier checks $|J|\le k$, checks for every selected subset $A_j$ and every $u\in A_j$ that
\[
g(u)\ge t,
\]
and then checks that the selected subsets cover all of $U$. This is again polynomial-time. Hence both problems lie in $\mathrm{NP}$. \qed

\subsection{Optimum-Preserving Reduction to \textsc{Exact-VAC-Dict}}

Fix an arbitrary \textsc{Set Cover} instance
\[
(\mathcal U,\mathcal S,k),
\qquad
\mathcal U=\{e_1,\dots,e_N\},
\qquad
\mathcal S=\{S_1,\dots,S_m\}.
\]
We construct an instance of \textsc{Exact-VAC-Dict} in polynomial time as follows.
\begin{enumerate}
    \item Set the ambient dimension to $n=1$.
    \item Encode each universe element $e_i$ by the rational point
    \[
    u_i:=i\in\mathbb Q.
    \]
    Thus
    \[
    U:=\{u_1,\dots,u_N\}\subseteq\mathbb Q.
    \]
    \item Define the target values by the constant table
    \[
    g(u_i):=0
    \qquad (1\le i\le N).
    \]
    \item For every set $S_j\in\mathcal S$, create one exact dictionary entry
    \[
    \eta_j=(A_j,a_j),
    \qquad
    A_j:=\{u_i: e_i\in S_j\},
    \qquad
    a_j(x):=0.
    \]
    \item Keep the same budget $k$.
\end{enumerate}
Call the resulting exact-side instance
\[
I_{=}(U,g,\mathcal D_{=},k).
\]

\paragraph{Lemma G.2 (All constructed exact entries are valid).}
Every dictionary entry $\eta_j=(A_j,a_j)$ in $I_{=}$ is valid.

\paragraph{Proof.}
For every $u\in A_j$ we have, by construction,
\[
g(u)=0=a_j(u).
\]
Thus each entry satisfies the exact-side validity condition. \qed

\paragraph{Lemma G.3 (YES/NO equivalence for the exact-side reduction).}
The original set-cover instance $(\mathcal U,\mathcal S,k)$ is a YES instance if and only if the constructed exact-side instance $I_{=}$ is a YES instance.

\paragraph{Proof.}
Assume first that $(\mathcal U,\mathcal S,k)$ is a YES instance. Then some index set
\[
J\subseteq \{1,\dots,m\},
\qquad |J|\le k,
\]
satisfies
\[
\mathcal U\subseteq \bigcup_{j\in J} S_j.
\]
By Lemma~G.2 all selected exact entries are valid. To show exact-side coverage, take any encoded point $u_i\in U$. Since the corresponding element $e_i$ lies in some selected set $S_j$, the point $u_i$ lies in the corresponding subset $A_j$. Hence
\[
U\subseteq \bigcup_{j\in J} A_j,
\]
so $I_{=}$ is a YES instance.

Conversely, suppose $I_{=}$ is a YES instance. Then some index set
\[
J\subseteq \{1,\dots,m\},
\qquad |J|\le k,
\]
selects valid entries and satisfies
\[
U\subseteq \bigcup_{j\in J} A_j.
\]
Take any universe element $e_i\in\mathcal U$. Its encoded point $u_i$ belongs to $U$, so coverage implies that $u_i\in A_j$ for some selected $j\in J$. By the definition of $A_j$, this is equivalent to $e_i\in S_j$. Hence every universe element is covered by the selected sets, so $(\mathcal U,\mathcal S,k)$ is also a YES instance. \qed

\paragraph{Lemma G.4 (Optimum preservation on the exact side).}
For the reduction above,
\[
\operatorname{OPT}_{=}(U,\mathcal D_{=})
=
\operatorname{OPT}_{\mathrm{SC}}(\mathcal U,\mathcal S).
\]

\paragraph{Proof.}
The reduction preserves the chosen index set exactly: selecting the subfamily $\{S_j:j\in J\}$ in the set-cover instance corresponds to selecting the exact dictionary entries indexed by the same $J$, and vice versa. Lemma~G.3 shows that cover feasibility is preserved in both directions, and the cardinality of the selected family is unchanged. Therefore the optimal values are equal. \qed

\subsection{Optimum-Preserving Reduction to \textsc{Thresh-VAC-Dict}}

We now use the same set-cover instance
\[
(\mathcal U,\mathcal S,k),
\qquad
\mathcal U=\{e_1,\dots,e_N\},
\qquad
\mathcal S=\{S_1,\dots,S_m\},
\]
and build a threshold-side dictionary instance.
\begin{enumerate}
    \item Reuse the rational point encoding
    \[
    U=\{u_1,\dots,u_N\}\subseteq\mathbb Q,
    \qquad
    u_i:=i.
    \]
    \item Reuse the constant value table
    \[
    g(u_i):=0
    \qquad (1\le i\le N).
    \]
    \item Fix the threshold
    \[
    t:=0.
    \]
    \item For each set $S_j\in\mathcal S$, create one threshold dictionary entry
    \[
    A_j:=\{u_i: e_i\in S_j\}.
    \]
    \item Keep the same budget $k$.
\end{enumerate}
Call the resulting threshold-side instance
\[
I_{\ge}=(U,g,t,\mathcal D_{\ge},k).
\]

\paragraph{Lemma G.5 (All constructed threshold entries are valid).}
Every threshold dictionary entry $A_j$ in $I_{\ge}$ is valid.

\paragraph{Proof.}
For every $u\in A_j$ we have, by construction,
\[
g(u)=0\ge 0=t.
\]
Hence each threshold entry is valid. \qed

\paragraph{Lemma G.6 (YES/NO equivalence for the threshold-side reduction).}
The original set-cover instance $(\mathcal U,\mathcal S,k)$ is a YES instance if and only if the constructed threshold-side instance $I_{\ge}$ is a YES instance.

\paragraph{Proof.}
If the selected sets indexed by $J$ cover $\mathcal U$, then the corresponding subsets $A_j$ cover $U$ by the same point-encoding argument as in Lemma~G.3, and all selected threshold entries are valid by Lemma~G.5. Hence $I_{\ge}$ is a YES instance.

Conversely, if the selected threshold entries indexed by $J$ cover $U$, then for every encoded point $u_i$ there exists some selected $A_j$ containing it. By construction this means $e_i\in S_j$, so the corresponding selected sets cover all of $\mathcal U$. Thus the original set-cover instance is a YES instance. \qed

\paragraph{Lemma G.7 (Optimum preservation on the threshold side).}
For the reduction above,
\[
\operatorname{OPT}_{\ge}(U,\mathcal D_{\ge},0)
=
\operatorname{OPT}_{\mathrm{SC}}(\mathcal U,\mathcal S).
\]

\paragraph{Proof.}
Exactly as on the exact side, the reduction preserves the chosen index set and its size. Lemma~G.6 gives feasibility preservation in both directions. Hence the optimum values agree. \qed

\subsection{NP-Completeness and Approximation Barrier}

The full decision theorem from Section~\ref{sec:vac-minimization} is now immediate.

\paragraph{Theorem G.8 (NP-completeness of explicit-dictionary selection).}
The decision problems \textsc{Exact-VAC-Dict} and \textsc{Thresh-VAC-Dict} are NP-complete.

\paragraph{Proof.}
By Proposition~G.1, both problems lie in $\mathrm{NP}$. Lemma~G.3 gives a polynomial-time reduction from \textsc{Set Cover} to \textsc{Exact-VAC-Dict}, and Lemma~G.6 gives a polynomial-time reduction from \textsc{Set Cover} to \textsc{Thresh-VAC-Dict}. Hence both dictionary-restricted decision problems are NP-hard. Combining NP-hardness with Proposition~G.1 yields NP-completeness. \qed

The structure of the reduction matters. The target function is constant, the threshold is constant, and every local dictionary entry is automatically valid. Therefore the hardness does \emph{not} come from local symbolic verification. It comes entirely from the combinatorics of selecting a smallest accepted subcover.

\paragraph{Lemma G.9 (Approximation-preserving transfer).}
Let $\rho:\mathbb N\to[1,\infty)$ be any function.
\begin{enumerate}
    \item If there exists a polynomial-time $\rho(N)$-approximation algorithm for \textsc{Exact-VAC-Dict}, then there exists a polynomial-time $\rho(N)$-approximation algorithm for \textsc{Set Cover}.
    \item If there exists a polynomial-time $\rho(N)$-approximation algorithm for \textsc{Thresh-VAC-Dict}, then there exists a polynomial-time $\rho(N)$-approximation algorithm for \textsc{Set Cover}.
\end{enumerate}
Here $N$ denotes the size of the universe of the original set-cover instance, equivalently the size of the encoded point domain $U$.

\paragraph{Proof.}
We prove the exact-side statement; the threshold-side proof is identical.

Assume there is a polynomial-time algorithm $\mathcal A_{=}$ that, on every exact dictionary instance with $|U|=N$, returns a valid subcover of size at most
\[
\rho(N)\,\operatorname{OPT}_{=}(U,\mathcal D_{=}).
\]
Given a set-cover instance $(\mathcal U,\mathcal S)$ with $|\mathcal U|=N$, form the exact-side dictionary instance $I_{=}$ by the reduction above and run $\mathcal A_{=}$ on it. Let $J$ be the returned index set. Mapping $J$ back to the set family
\[
\{S_j:j\in J\}\subseteq\mathcal S
\]
produces a set cover of $\mathcal U$ by Lemma~G.3. Moreover,
\[
|J|
\le
\rho(N)\,\operatorname{OPT}_{=}(U,\mathcal D_{=})
=
\rho(N)\,\operatorname{OPT}_{\mathrm{SC}}(\mathcal U,\mathcal S),
\]
where the equality is Lemma~G.4. Hence we obtain a polynomial-time $\rho(N)$-approximation for \textsc{Set Cover}. The threshold-side statement follows from the same argument using Lemma~G.7. \qed

\paragraph{Corollary G.10 (Logarithmic approximation barrier for explicit-dictionary minimization).}
For every fixed $\varepsilon\in(0,1)$, unless $\mathrm P=\mathrm{NP}$, there is no polynomial-time algorithm that approximates either
\[
\operatorname{OPT}_{=}(U,\mathcal D_{=})
\qquad\text{or}\qquad
\operatorname{OPT}_{\ge}(U,\mathcal D_{\ge},t)
\]
within factor
\[
(1-\varepsilon)\ln N,
\qquad N:=|U|.
\]

\paragraph{Proof.}
The classical logarithmic approximation barrier for \textsc{Set Cover} goes back to Feige~\cite{Feige1998SetCover}, and the fixed-$\varepsilon$ hardness form used in Section~\ref{sec:vac-minimization} is given by Dinur and Steurer~\cite{DinurSteurer2014}. If either associated dictionary-restricted optimization problem admitted a polynomial-time $(1-\varepsilon)\ln N$ approximation, Lemma~G.9 would transfer the same guarantee to \textsc{Set Cover}, contradicting those lower bounds. \qed

This is the natural approximation scale for classical cover minimization within the unrestricted explicit-dictionary model. Thus the selection layer is not a benign post-processing step after candidate generation; even after candidate cells are explicit, it already inherits the standard logarithmic barrier of set-family optimization.

\subsection{Boundary Cases and Scope Clarifications}

We close by recording exactly what the appendix establishes and what it does not.

\paragraph{Proposition G.11 (Threshold-side hardness already holds on the true-claim subclass).}
Restrict \textsc{Thresh-VAC-Dict} to instances satisfying
\[
\inf_{u\in U} g(u)\ge t.
\]
The restricted decision and optimization problems remain NP-hard, and the optimum-preserving reduction above still applies.

\paragraph{Proof.}
In every reduced threshold instance $I_{\ge}$ we set
\[
g\equiv 0
\qquad\text{on }U,
\qquad
 t=0.
\]
Hence
\[
\inf_{u\in U} g(u)=0=t.
\]
So the reduction lands entirely inside the subclass where the global threshold claim is true. Lemmas~G.6--G.7 therefore already prove NP-hardness and optimum preservation on that restricted subclass. \qed

This clarifies an important boundary case. On the threshold side, false claims automatically force semantic threshold VAC to be $+\infty$ by Proposition~F.12 of Appendix~\ref{app:vac-bounds}. The hardness result here is not exploiting that trivial obstruction; it survives on the much smaller subclass where the threshold claim itself is globally true.

\paragraph{Proposition G.12 (Trivial conventions and redundant entries).}
For both dictionary problems:
\begin{enumerate}
    \item if $U=\varnothing$, then the optimum value is $0$;
    \item deleting an invalid dictionary entry does not change the optimum value;
    \item deleting an empty dictionary entry does not increase the optimum value;
    \item deleting a duplicate copy of an already-present valid entry does not increase the optimum value.
\end{enumerate}

\paragraph{Proof.}
If $U=\varnothing$, the empty subfamily already covers the domain, so the optimum is $0$. An invalid entry can never appear in a feasible solution by definition, hence deleting it does not change feasibility or the optimum. An empty entry covers no points, so deleting it cannot remove any feasible cover. A duplicate valid entry contributes the same subset as an existing one, so retaining one copy is always as good as retaining both. \qed

Thus the NP-hardness statement is not an artifact of pathological dictionary encodings. It persists after the obvious syntactic cleanup operations.

\paragraph{(i) The hardness is dictionary-restricted.}
The input already contains the full finite candidate pool. No claim is made here about the unrestricted decision complexity of semantic VAC over arbitrary infinite continuous languages.

\paragraph{(ii) The hardness is purely combinatorial.}
In the constructed instances, local validity is trivial:
\[
g\equiv 0,
\qquad
 a_j\equiv 0,
\qquad
 t=0.
\]
Hence the obstruction comes entirely from selecting a smallest accepted subcover, not from generating cells or proving local soundness.

\paragraph{(iii) What the reduction proves unconditionally.}
The reduction proves hardness for the unrestricted explicit-dictionary selection problem: the input already contains the full finite candidate family, and the only remaining task is to choose a smallest accepted subcover.

\paragraph{Definition G.13 (Pipeline-exposed dictionary class).}
Fix a family $\Pi$ of proof-producing mechanisms. For each verification instance $I$, let $\mathcal D_\Pi(I)$ denote the finite dictionary of verifier-visible candidate cells exposed by $\Pi$ before any cover-minimization step. The exposed dictionary class of $\Pi$ is
\[
\mathfrak C_\Pi:=\{\mathcal D_\Pi(I): I\in\mathcal I_\Pi\},
\]
where $\mathcal I_\Pi$ is the instance family on which $\Pi$ is defined.

\paragraph{Definition G.14 (Set-cover-rich exposed family class).}
We say that $\Pi$ is \emph{set-cover-rich} if there is a polynomial-time mapping from every set-cover instance $(\mathcal U,\mathcal S,k)$ to a verification instance $I$ such that feasible accepted subcovers of $\mathcal D_\Pi(I)$ correspond to set covers of $(\mathcal U,\mathcal S)$. If the correspondence preserves optimum cardinality, we say that the realization is \emph{optimum-preserving}.

\paragraph{Theorem G.15 (Conditional transfer to concrete pipelines).}
Let $\Pi$ be a proof-producing mechanism whose exposed dictionary class is set-cover-rich. Then the induced selection problem over dictionaries in $\mathfrak C_\Pi$ is NP-hard. If the realizing map is optimum-preserving, then the logarithmic inapproximability barrier of Corollary~G.10 transfers as well.

\paragraph{Proof sketch.}
Compose the classical set-cover reduction used above with the realizing map into $\mathfrak C_\Pi$. By hypothesis, feasible accepted subcovers correspond to set covers, and in the optimum-preserving case the optimum values coincide. NP-hardness and logarithmic inapproximability therefore transfer directly. \qed

\paragraph{Remark G.16 (Why the richness hypothesis cannot be dropped).}
Without a representability hypothesis, no generic transfer from unrestricted explicit dictionaries to arbitrary pipeline-generated candidate pools is valid. A concrete pipeline may expose only a highly structured class of candidate families, and within such classes minimum-cover problems can be substantially easier than general set cover. The present appendix therefore proves unconditional hardness only for unrestricted explicit dictionaries. Characterizing which pipeline-generated family classes are set-cover-rich, and which admit polynomial-time minimization, is a natural open problem beyond the scope of the present paper.

\paragraph{(iv) Relation to the proof-language layer.}
Section~\ref{sec:vac-proof-languages} shows that, once a concrete syntax realizes semantic covers with only constant-factor leaf overhead, minimum leaf count tracks the corresponding semantic selection problem up to that fixed realization factor. Combined with the results above, this identifies a genuine combinatorial complexity source at the explicit-dictionary selection layer. Transfer to any particular proof pipeline, however, requires that the exposed candidate-family class satisfy an additional representability hypothesis such as Theorem~G.15.

In that precise sense, Section~\ref{sec:vac-minimization} is about \emph{selection-layer hardness}, and this appendix supplies the full reductions behind that claim.
